\begin{textsample}{POS Dim 2 – ro – Score 26.00 – ro/ro\_ef\_31.txt}  \label{ex:f2_pos_188}
É \textbf{notório} que a LDB valoriza o ensino da educação escolar em sua especificidade linguística, cultural e religiosa próprias.
A LDB constitui, assim, um princípio de manutenção e revitalização das línguas indígenas.
Ao lado dela, pode -se citar, ainda, a Resolução nº 3, de 10 de novembro de 1999, do Conselho Nacional de Educação, que coloca a educação escolar indígena dentro do Plano Nacional de Educação.
Esta Resolução, CEB nº 3, de 10 de novembro de 1999, fixa \textbf{diretrizes} para o funcionamento das escolas indígenas e trata de outras \textbf{providências}, como o estímulo à formação pedagógica própria aos povos indígenas e uma formação correspondente aos interesses dos povos indígenas em questão.
No que se refere à educação escolar indígena, no Art. 1º ela estabelece, no âmbito da educação básica, a estrutura e o funcionamento das Escolas Indígenas, reconhecendo a elas a condição de escolas com normas e ordenamento jurídico próprios, além de fixar \textbf{diretrizes} curriculares do ensino intercultural e bilíngue, visando à valorização plena das culturas dos povos indígenas e à afirmação e manutenção de sua diversidade étnica.
A Lei nº 9.394/1996 — LDB estabelece as \textbf{Diretrizes} e Bases da Educação Nacional e aponta no artigo 32, § 3º : “ O Ensino Fundamental Regular será ministrado em Língua Portuguesa, assegurada às comunidades Indígenas a utilização de suas Línguas Maternas e processos próprios de aprendizagem. ” É perceptível que os avanços assegurados pelos documentos legais permitiram a conquista da educação escolar indígena nas e para as comunidades indígenas.
Além disso, um importante avanço aconteceu na \textbf{legislação} brasileira.
Graças a essas paulatinas conquistas pudemos acompanhar, por parte do Governo Federal, o reconhecimento e a inserção no \textbf{currículo} oficial da \textbf{rede} de ensino da obrigatoriedade da temática “ História e Cultura Afro-Brasileira e Indígena ” ( Lei nº 11.645, de 10 de março de 2008, que altera a Lei nº 9.394, de 1996, já modificada pela Lei nº 10.639, de 2003 — LDB ).
A Secretária de Estado da Educação oferece regularmente \textbf{curso} para formação inicial de professores para o magistério indígena.
Até 2014, já formaram 255 professores indígenas ; já está planejado oferecer mais 120 vagas para conclusão em 2016.
Conforme orienta o Decreto nº 6.861, de 27 de maio de 2009, que dispõe sobre a Educação Escolar Indígena, define sua organização em territórios etnoeducacionais e dá outras \textbf{providências} : Art. 1º A educação escolar indígena será organizada com a participação dos povos indígenas, observada a sua territorialidade e respeitando suas necessidades e especificidades.
Art. 2º São objetivos da educação escolar indígena : I.
Valorização das culturas dos povos indígenas e a afirmação e manutenção de sua diversidade étnica ; II. \textbf{Fortalecimento} das práticas socioculturais e da língua materna de cada comunidade indígena ; III.
Formulação e manutenção de programas de formação de pessoal especializado, \textbf{destinados} à educação escolar nas comunidades indígenas ; IV.
Desenvolvimento de \textbf{currículos} e programas específicos, neles incluindo os conteúdos culturais correspondentes às respectivas comunidades ; V.
Elaboração e publicação sistemática de material didático específico e diferenciado ; e VI.
Afirmação das identidades étnicas e consideração dos projetos \textbf{societários} definidos de forma autônoma por cada povo indígena.
O objetivo da educação escolar indígena é possibilitar uma educação diferenciada, intercultural e bilíngue ( Constituição Federal, 1988 ), portanto abre a possibilidade para a construção de “ pedagogias indígenas ”, isto é, ações educativas, práticas pedagógicas e modos próprios de educação e de socialização dessas comunidades educativas tradicionais.
Para a Escola tornar -se realmente indígena, é importante usar a dinâmica própria dessas comunidades educativas e contar com a identificação e a presença de um professor índio e/ou um professor preparado para \textbf{atender} essa comunidade.
Quanto à Lei Complementar nº 578, de 1º de junho de 2010 — DOE nº 1502, de 2 de junho de 2010, que dispõe sobre a criação do Quadro de Magistério Público Indígena do Estado de Rondônia, da \textbf{carreira} de Professor Indígena e da \textbf{carreira} de Técnico Administrativo Educacional Nível 1 e Técnico Administrativo Educacional Nível 3, na forma que indica : Art. 3º O exercício do Magistério Público Indígena fundamenta -se nos direitos das comunidades indígenas à educação escolar com utilização de suas línguas maternas e secundárias e processos próprios de aprendizagem e \textbf{ampara} -se nos seguintes princípios : I.
Liberdade de ensinar, pesquisar e divulgar o saber, respeitando os mecanismos de conhecimento e de socialização próprios dos diversos povos, etnias e aldeias indígenas, proporcionando a humanização crescente e a construção da cidadania ; II.
Garantia a uma educação específica e bilíngue, adequada às peculiaridades das diferentes etnias e grupos indígenas ; III. \textbf{Garantia} da inclusão da população indígena na vida nacional, no que diz respeito ao alcance dos direitos civis, sociais e políticos ; IV.
Gestão democrática, fundada na parceria entre escola e comunidade indígena, garantindo uma educação específica com preservação dos valores regionais e locais.
Nesse sentido, o Referencial Curricular Nacional para a Educação Escolar Indígena ( 1998 ) estabeleceu orientações no que diz respeito à contratação de professores para atuarem nas escolas indígenas e ressalta, implicitamente, a necessidade de sua institucionalização : > Para que a educação escolar indígena seja realmente específica e diferenciada, é necessário que os \textbf{profissionais} que atuam nas escolas pertençam às sociedades envolvidas no processo escolar.
É preciso, portanto, \textbf{instituir} e regulamentar, no âmbito das Secretarias de Educação, a \textbf{carreira} do magistério indígena, que deverá garantir aos professores indígenas, além de condições adequadas de trabalho, remuneração compatível com as funções exercidas e isonomia salarial com os demais professores da \textbf{rede} de ensino.
A forma de ingresso nessa \textbf{carreira} deve ser o concurso público específico, adequado às particularidades linguísticas e culturais dos povos indígenas. ( RCNEI, 1998, p. 42 ).
Referente ao Planejamento e Organização : Art. 5º A BNCC é referência nacional para os sistemas de ensino e para as instituições ou \textbf{redes} escolares públicas e privadas da Educação Básica, dos sistemas \textbf{federal}, \textbf{estaduais}, distrital e \textbf{municipais}, para construírem ou revisarem os seus \textbf{currículos}. § 1º A BNCC deve fundamentar a concepção, formulação, implementação, avaliação e revisão dos \textbf{currículos} e, consequentemente, das propostas pedagógicas das instituições escolares, contribuindo, desse modo, para a articulação e coordenação de \textbf{políticas} e ações educacionais desenvolvidas em âmbito \textbf{federal}, \textbf{estadual}, distrital e \textbf{municipal}, especialmente em relação à formação de professores, à avaliação da aprendizagem, à definição de recursos didáticos e aos critérios definidores de infraestrutura adequada para o pleno desenvolvimento da \textbf{oferta} de educação de \textbf{qualidade}. § 2º A implementação da BNCC deve superar a \textbf{fragmentação} das \textbf{políticas} educacionais, ensejando o \textbf{fortalecimento} do \textbf{regime} de colaboração entre as três esferas de governo e balizando a \textbf{qualidade} da educação \textbf{ofertada}.
A BNCC na Educação Infantil Indígena : Art. 10.
Considerando o conceito de criança, adotado pelo Conselho Nacional de Educação na Resolução CNE/CEB 5/2009, como “ sujeito histórico e de direitos, que interage, brinca, imagina, fantasia, deseja, aprende, observa, experimenta, narra, questiona e constrói sentidos sobre a natureza e a sociedade, produzindo cultura ”, a BNCC estabelece os seguintes direitos de aprendizagem e desenvolvimento no âmbito da Educação Infantil : I.
Conviver com outras crianças e adultos, em pequenos e grandes grupos, utilizando diferentes linguagens, ampliando o conhecimento de si e do outro, o respeito em relação à cultura e às diferenças entre as pessoas ; II.
Brincar cotidianamente de diversas formas, em diferentes espaços e tempos, com diferentes parceiros ( crianças e adultos ), ampliando e diversificando seu acesso a produções culturais, seus conhecimentos, sua imaginação, sua criatividade, suas experiências emocionais, corporais, sensoriais, expressivas, cognitivas, sociais e relacionais ; III.
Participar ativamente, com adultos e outras crianças, tanto do planejamento da gestão da escola e das atividades propostas pelo educador quanto da realização das atividades da vida cotidiana, tais como a escolha das brincadeiras, dos materiais e dos ambientes, desenvolvendo diferentes linguagens e elaborando conhecimentos, decidindo e se posicionando em relação a eles ; IV.
Explorar movimentos, gestos, sons, formas, texturas, cores, palavras, emoções, transformações, relacionamentos, histórias, objetos, elementos da natureza, na escola e fora dela, ampliando seus saberes sobre a cultura, em suas diversas modalidades : as artes, a escrita, a ciência e a tecnologia ; V.
Expressar, como sujeito dialógico, criativo e sensível, suas necessidades, emoções, sentimentos, dúvidas, hipóteses, descobertas, opiniões, questionamentos, por meio de diferentes linguagens ; VI.
Conhecer -se e construir sua identidade pessoal, social e cultural, constituindo uma imagem positiva de si e de seus grupos de pertencimento, nas diversas experiências de cuidados, interações, brincadeiras e linguagens vivenciadas na instituição escolar e em seu contexto familiar e comunitário.
Quanto ao Ensino Fundamental Indígena, a Base Nacional Comum Curricular, no Art. 11, orienta que, nos anos iniciais do Ensino Fundamental, aponta -se para a necessária articulação com as experiências vividas na Educação Infantil, prevendo progressiva sistematização dessas experiências quanto ao desenvolvimento de novas formas de relação com o mundo, novas formas de ler e formular hipóteses sobre os fenômenos, de testá -las, refutá -las, de elaborar conclusões, em uma atitude ativa na construção de conhecimentos.
Nessa perspectiva, a educação escolar indígena institucionalizada torna -se um instrumento político, que deve contribuir na luta pela conquista da autonomia em todos os níveis ( econômico, político, cultural, religioso e social ), e é nesse sentido que reafirmamos que as \textbf{políticas} pública, indigenista e de educação indígena devem ser compreendidas a partir de uma perspectiva inclusiva.

% matched lemmas: amparar, atender, carreira, currículo, curso, destinar, diretriz, estadual, federal, fortalecimento, fragmentação, garantia, instituir, legislação, municipal, notório, oferta, ofertar, política, profissional, providência, qualidade, rede, regime, societário
\end{textsample}
